\section*{Capítulo \ref{ch:g10:concentration-and-dilution} --- Concentración y dilución}

\begin{solution}{exo:g10:concentration-and-dilution:1}
$C_m = \qty{5.0}{g} / \qty{0.250}{L} = \qty{20}{g/L}$.
\end{solution}

\begin{solution}{exo:g10:concentration-and-dilution:2}
$c = \qty{0.050}{mol} / \qty{0.200}{L} = \qty{0.25}{mol/L}$.
\end{solution}

\begin{solution}{exo:g10:concentration-and-dilution:3}
$M(\ce{CuSO4}) = 63.5 + 32.1 + 4 \times 16.0 = \qty{159.6}{g/mol}$;
$n = 0.10 \times 0.1000 = \qty{0.0100}{mol}$; $m = 0.0100 \times 159.6 =
\qty{1.60}{g}$.
\end{solution}

\begin{solution}{exo:g10:concentration-and-dilution:4}
$F = 1.0 / 0.050 = 20$; $V_0 = \qty{100.0}{mL} / 20 = \qty{5.0}{mL}$.
\end{solution}

\begin{solution}{exo:g10:concentration-and-dilution:5}
$C_m = c \times M = 0.10 \times 58.5 = \qty{5.85}{g/L}$, unos
\qty{5.9}{g/L}.
\end{solution}

\begin{solution}{exo:g10:concentration-and-dilution:6}
$V_0 = c_1 V_1 / c_0 = 0.020 \times 250.0 / 0.50 = \qty{10.0}{mL}$
(\omterm{def:g10:concentration-and-dilution:dilution}{factor de dilución} 25). Una pipeta aforada de \qty{10.0}{mL} y un
matraz aforado de \qty{250.0}{mL}.
\end{solution}

\begin{solution}{exo:g10:concentration-and-dilution:7}
El matraz aforado tiene un único enrase, hecho para un solo volumen y
exacto a una pequeña fracción de un tanto por ciento; las graduaciones
de un vaso de precipitados son solo orientativas. Del mismo modo, una
pipeta aforada da un volumen con mucha exactitud, mientras que una
probeta es más ancha y se lee con menos precisión.
\end{solution}

\begin{solution}{exo:g10:concentration-and-dilution:8}
Entre \qty{0.02}{mol/L} y \qty{0.04}{mol/L}: alrededor de
\qty{0.03}{mol/L}. Para más precisión, se añaden patrones entre esos dos,
por ejemplo de 0.025, 0.030 y \qty{0.035}{mol/L}: una escala más fina
allí donde hace falta.
\end{solution}

\begin{solution}{exo:g10:concentration-and-dilution:9}
$M(\ce{C6H12O6}) = \qty{180.0}{g/mol}$; $c = 50 / 180.0 =
\qty{0.28}{mol/L}$.
\end{solution}

\begin{solution}{exo:g10:concentration-and-dilution:10}
$M = \qty{342.0}{g/mol}$, así que $n = \qty{1.00}{mol}$ y $c =
\qty{1.00}{mol/L}$. Diluido diez veces: $c = \qty{0.100}{mol/L}$, es
decir, $C_m = 0.100 \times 342.0 = \qty{34.2}{g/L}$. Una cucharada de
\qty{0.020}{L} contiene $34.2 \times 0.020 = \qty{0.68}{g}$ de sacarosa.
\end{solution}

\begin{solution}{exo:g10:concentration-and-dilution:11}
(a) La misma masa de \omterm{def:g6:solutions-and-solubility:solute}{soluto} en el doble de volumen: \qty{6.0}{g/L}.
(b) El volumen se reduce más o menos a la mitad y la masa de \omterm{def:g6:solutions-and-solubility:solute}{soluto} no
cambia: unos \qty{24}{g/L}.
\end{solution}

\begin{solution}{exo:g10:concentration-and-dilution:12}
\begin{enumerate}
  \item Tres \omterm{def:g10:concentration-and-dilution:dilution}{diluciones} por 10 equivalen a una \omterm{def:g10:concentration-and-dilution:dilution}{dilución} por $10^3$:
        $\qty{0.10}{mol/L} / 1000 = \qty{1.0e-4}{mol/L}$.
  \item Una sola \omterm{def:g10:concentration-and-dilution:dilution}{dilución} por 1000 exigiría o bien un volumen diminuto
        de \omterm{def:g10:concentration-and-dilution:dilution}{disolución madre} (\qty{0.1}{mL} en \qty{100}{mL}), imposible
        de medir con exactitud con una pipeta, o bien un matraz muy
        grande (\qty{1}{mL} en \qty{1}{L}). Tres \omterm{def:g10:concentration-and-dilution:dilution}{diluciones} por diez usan
        pipetas y matraces corrientes.
\end{enumerate}
\end{solution}

\begin{solution}{exo:g10:concentration-and-dilution:13}
\begin{enumerate}
  \item $M(\ce{CaCl2}) = 40.1 + 2 \times 35.5 = \qty{111.1}{g/mol}$;
        $n = 11.1 / 111.1 = \qty{0.100}{mol}$; $c = 0.100 / 0.5000 =
        \qty{0.200}{mol/L}$.
  \item Cada \ce{CaCl2} da un \ce{Ca^2+} y dos \ce{Cl-}:
        $[\ce{Ca^2+}] = \qty{0.200}{mol/L}$ y $[\ce{Cl-}] =
        \qty{0.400}{mol/L}$.
\end{enumerate}
\end{solution}

\begin{solution}{exo:g10:concentration-and-dilution:14}
\begin{enumerate}
  \item $r = \qty{0.60}{cm}$, $h = \qty{0.20}{cm}$:
        $V = \pi \times 0.60^2 \times 0.20 = \qty{0.23}{cm^3}$, es
        decir, \qty{0.23}{mL} de agua de más.
  \item El mismo \omterm{def:g6:solutions-and-solubility:solute}{soluto} está en \qty{100.23}{mL} en lugar de
        \qty{100.00}{mL}: la concentración es demasiado baja en
        $0.23 / 100.23 \approx
        \qty{0.23}{\percent}$.
\end{enumerate}
\end{solution}

\begin{solution}{exo:g10:concentration-and-dilution:15}
$n = 0.20 \times 0.100 + 0.10 \times 0.300 = 0.020 + 0.030 =
\qty{0.050}{mol}$ en \qty{0.400}{L}: $c = \qty{0.125}{mol/L}$. No es la
media simple, \qty{0.15}{mol/L}: hay tres veces más de la disolución
diluida, así que la \omterm{def:g6:pure-substances-and-mixtures:mixture}{mezcla} está más cerca de \qty{0.10}{mol/L}. Es la
media ponderada por los volúmenes.
\end{solution}

\begin{solution}{pb:g10:concentration-and-dilution:1}
\textbf{1.} \qty{0.9}{g} por \qty{0.100}{L}: $C_m = \qty{9.0}{g/L}$.

\textbf{2.} \qty{9.0}{g} por \qty{1000}{mL}, es decir, \qty{9.0}{mg} por
mililitro.

\textbf{3.} $m = 9.0 \times 0.500 = \qty{4.5}{g}$.

\textbf{4.} Sí: unos \qty{9}{g} por litro frente a unos \qty{360}{g} por
kilogramo de agua en la saturación, unas cuarenta veces menos.

\textbf{5.} $9.0 \times 1.0 = \qty{9.0}{g}$ de sal.

\textbf{6.} $M(\ce{NaCl}) = 23.0 + 35.5 = \qty{58.5}{g/mol}$.

\textbf{7.} $c = 9.0 / 58.5 = \qty{0.154}{mol/L}$.

\textbf{8.} $n = 0.154 \times 0.500 = \qty{0.0769}{mol}$ (o
$4.5 / 58.5$).

\textbf{9.} Cada \ce{NaCl} da un \ce{Na+} y un \ce{Cl-}: los dos a
\qty{0.154}{mol/L}.

\textbf{10.} $0.0769 \times \num{6.02e23} = \num{4.63e22}$ \omterm{def:g9:ions:ion}{iones} sodio.

\textbf{11.} $F = 200 / 9.0 = 22.2$.

\textbf{12.} Los \qty{4.5}{g} de la bolsa vienen todos de la disolución
concentrada: $V_0 = \qty{4.5}{g} / \qty{200}{g/L} = \qty{0.0225}{L} =
\qty{22.5}{mL}$.

\textbf{13.} $V_0 = V_1 / F = 500 / 22.2 = \qty{22.5}{mL}$.

\textbf{14.} $c_0 = 200 / 58.5 = \qty{3.42}{mol/L}$;
$c_0 V_0 = 3.42 \times 0.0225 = \qty{0.0769}{mol}$ y
$c_1 V_1 = 0.154 \times 0.500 = \qty{0.0769}{mol}$: iguales.

\textbf{15.} Unos $500 - 22.5 = \qty{478}{mL}$ de agua estéril.

\textbf{16.} Una pipeta graduada de \qty{25}{mL} (o una bureta), leída
a \qty{0.1}{mL}.

\textbf{17.} En un matraz aforado de \qty{500.0}{mL}, completado hasta
el enrase.

\textbf{18.} $C_m = 200 \times 23.0 / 500 = \qty{9.20}{g/L}$, demasiado
concentrada en $0.20 / 9.0 \approx \qty{2.2}{\percent}$.

\textbf{19.} \qty{22.5}{mL} de la disolución concentrada (``20 \%'')
para una bolsa de \qty{500}{mL}.
\end{solution}
